\chapter{ऐल्कोहॉल: OH समूह का सक्रियण}\label{ch:b1:alcohol-activation}

ऐल्कोहॉल हर जगह हैं: ऐल्कीनों से और किण्वन से टनों में बनाए जाते हैं,
और पिछले अध्याय का हर ग्रीन्यार योग इन्हें देता है। फिर भी कोई ऐल्कोहॉल
हैलोऐल्केन की तरह प्रतिस्थापन से नहीं गुज़रता: OH समूह को हाइड्रॉक्साइड आयन के रूप में
निकलना पड़ता, जो एक \omterm{def:b1:predominant-reaction:strong-acid}{प्रबल क्षारक} और अत्यंत खराब \omterm{def:b1:electronic-effects:nucleophile}{निर्गामी समूह} है (\cref{ch:b1:nucleophilic-substitution})।
रसायनज्ञों के पास तीन उत्तर हैं। इसके बजाय OH का प्रोटॉन हटा दीजिए, तो
\omterm{def:b1:alcohol-activation:alkoxide}{ऐल्कॉक्साइड} एक शक्तिशाली \omterm{def:b1:electronic-effects:nucleophile}{नाभिकरागी} बन जाता है। एक प्रोटॉन जोड़िए, तो \omterm{def:b1:electronic-effects:nucleophile}{निर्गामी
समूह} जल बन जाता है। या हाइड्रोजन को ऐसे समूह से बदलिए जो
ऑक्सीजन को एक उत्कृष्ट \omterm{def:b1:electronic-effects:nucleophile}{निर्गामी समूह} का भाग बना दे, कार्बन को छुए
बिना। यह अध्याय इन तीनों मार्गों का, और उनके साथ आने वाले
क्रियाविधि और त्रिविम रसायन के विकल्पों का अनुसरण करता है।

\begin{recall}
SN1, SN2 और उनका त्रिविम रसायन (\cref{ch:b1:nucleophilic-substitution});
E1, E2 और \omterm{prop:b1:elimination:zaitsev}{ज़ैत्सेव का नियम} (\cref{ch:b1:elimination}); \omterm{def:b1:predominant-reaction:acidity-constant}{अम्लता स्थिरांक}
और प्रधान-अभिक्रिया विधि (\cref{ch:b1:predominant-reaction});
ऐल्कोहॉलों की अम्लता पर ऐल्किल समूहों का प्रभाव
(\cref{exo:b1:electronic-effects:11})।
\end{recall}

\section{अम्ल के रूप में ऐल्कोहॉल: ऐल्कॉक्साइड}

\begin{definition}[ऐल्कॉक्साइड]\label{def:b1:alcohol-activation:alkoxide}
\emph{ऐल्कॉक्साइड}\index{ऐल्कॉक्साइड} वह संयुग्मी क्षारक \ce{RO-} है जो किसी
ऐल्कोहॉल \ce{ROH} से व्युत्पन्न होता है: मेथॉक्साइड \ce{CH3O-}, एथॉक्साइड \ce{CH3CH2O-},
टर्ट-ब्यूटॉक्साइड \ce{(CH3)3CO-}। ऐल्कॉक्साइड \omterm{def:b1:predominant-reaction:strong-acid}{प्रबल क्षारक} और अच्छे
\omterm{def:b1:electronic-effects:nucleophile}{नाभिकरागी} होते हैं।
\end{definition}

\begin{proposition}[ऐल्कोहॉलों की अम्लता]\label{prop:b1:alcohol-activation:alcohol-acidity}
ऐल्कोहॉल जल से दुर्बल या लगभग उतने ही \omterm{def:b1:predominant-reaction:strong-acid}{दुर्बल अम्ल} हैं: $\mathrm{p}K_a$
15.5 (मेथेनॉल), 15.9 (एथेनॉल), 17.1 (प्रोपेन-2-ऑल), 19.2
(2-मेथिलप्रोपेन-2-ऑल), जबकि युग्म \ce{H2O}/\ce{OH-} के लिए 14.0।
इसलिए हाइड्रॉक्साइड उनका केवल आंशिक विप्रोटॉनन करता है; \omterm{def:b1:alcohol-activation:alkoxide}{ऐल्कॉक्साइड} में
पूर्ण रूपांतरण के लिए अधिक \omterm{def:b1:predominant-reaction:strong-acid}{प्रबल क्षारक} या कोई क्षार धातु चाहिए।
\end{proposition}
% ledger: pka:methanol, pka:ethanol, pka:2-propanol, pka:tBuOH

\begin{proof}
अभिक्रिया \ce{ROH + OH- <=> RO- + H2O} का $K = 10^{14.0 -
\mathrm{p}K_a(\ce{ROH})}$ है, $10^{-1.5}$ और $10^{-5}$ के बीच: प्रतिकूल।
सोडियम धातु प्रोटॉन को अपचित करती है, \ce{2ROH + 2Na -> 2RO- + 2Na+ + H2}, और
सोडियम हाइड्राइड का हाइड्राइड आयन उसे अनुत्क्रमणीय रूप से ले लेता है,
\ce{ROH + H- -> RO- + H2}: हाइड्रोजन निकल जाती है, और अभिक्रिया
पूर्ण होती है, उसका साम्य स्थिरांक चाहे जो हो।
\end{proof}

\begin{inthelab}[सोडियम हाइड्राइड]
सोडियम हाइड्राइड खनिज तेल में परिक्षिप्त धूसर चूर्ण के रूप में बिकता है, जो
उसे वायु की नमी से बचाता है। यह जल से प्रचंड अभिक्रिया करता है,
ज्वलनशील हाइड्रोजन मुक्त करते हुए; उसे जल्दी तौला जाता है, अक्रिय गैस के अंतर्गत
शुष्क ऐल्कोहॉल में भागों में मिलाया जाता है, और हाइड्रोजन के बुलबुले
विप्रोटॉनन की प्रगति दिखाते हैं। अंत में आधिक्य हाइड्राइड को किसी ऐल्कोहॉल को
धीरे-धीरे मिलाकर नष्ट किया जाता है, कभी जल से नहीं।
\end{inthelab}

\section{विलियमसन ईथर संश्लेषण}

\begin{definition}[विलियमसन ईथर संश्लेषण]\label{def:b1:alcohol-activation:williamson}
\emph{विलियमसन ईथर संश्लेषण}\index{विलियमसन ईथर संश्लेषण}
ईथर $\mathrm{R{-}O{-}R'}$ बनाता है, \omterm{def:b1:alcohol-activation:alkoxide}{ऐल्कॉक्साइड} \ce{RO-} की
किसी हैलोऐल्केन (या \omterm{def:b1:alcohol-activation:sulfonate}{सल्फ़ोनेट एस्टर}) $\mathrm{R'{-}X}$ पर SN2 अभिक्रिया द्वारा।
\end{definition}

\begin{omfigure}
\setchemfig{atom sep=2.1em}
\schemestart
\chemfig{CH_3CH_2-@{o}O^{-}}\,\chemfig{Na^{+}}
\+
\chemfig{@{c}CH_3-[@{ci}]@{i}I}
\arrow{->}
\chemfig{CH_3CH_2-O-CH_3}
\+
\chemfig{Na^{+}}\,\chemfig{I^{-}}
\schemestop

\medskip
\chemmove{\draw[->, thick, shorten <=2pt, shorten >=2pt](o) .. controls +(80:8mm) and +(100:8mm) .. (c);
          \draw[->, thick, shorten <=2pt, shorten >=2pt](ci) .. controls +(90:5mm) and +(90:5mm) .. (i);}

{\small मेथॉक्सीएथेन का विलियमसन संश्लेषण: एथॉक्साइड आयोडोमेथेन के
कार्बन पर आयोडीन के विपरीत ओर से आक्रमण करता है (SN2)।}
\end{omfigure}

\begin{method}[दोनों भागीदारों का चयन]\label{met:b1:alcohol-activation:williamson-choice}
ईथर $\mathrm{R{-}O{-}R'}$ को ऑक्सीजन के किसी भी ओर से काटा जा सकता है।
\begin{enumerate}
  \item दोनों संभावित जोड़े लिखिए: \ce{RO-} के साथ $\mathrm{R'X}$, और
        $\mathrm{R'O^-}$ के साथ \ce{RX}।
  \item वह जोड़ा रखिए जिसमें हैलाइड मेथिल या प्राथमिक हो: SN2 को
        अबाधित कार्बन चाहिए।
  \item तृतीयक हैलाइड वाला जोड़ा अस्वीकार कीजिए: \omterm{def:b1:alcohol-activation:alkoxide}{ऐल्कॉक्साइड}, जो \omterm{def:b1:predominant-reaction:strong-acid}{प्रबल क्षारक} है,
        इसके बजाय विलोपन (E2) करेगा।
  \item द्वितीयक हैलाइड के साथ ईथर और ऐल्कीन के मिश्रण की अपेक्षा कीजिए।
\end{enumerate}
\end{method}

\begin{example}[एक असममित ईथर]\label{ex:b1:alcohol-activation:williamson}
2-मेथॉक्सी-2-मेथिलप्रोपेन (मेथिल टर्ट-ब्यूटिल ईथर): आयोडोमेथेन के साथ सोडियम टर्ट-ब्यूटॉक्साइड
इसे स्वच्छ रूप से देता है; तृतीयक
ब्रोमाइड, 2-ब्रोमो-2-मेथिल\-प्रोपेन, के साथ सोडियम मेथॉक्साइड E2 द्वारा 2-मेथिलप्रोपीन देता है। \omterm{def:b1:alcohol-activation:alkoxide}{ऐल्कॉक्साइड} की ओर
भारी-भरकम, हैलाइड की ओर छोटा।
\end{example}

\section{OH समूह का सक्रियण}

\begin{definition}[सल्फ़ोनेट एस्टर]\label{def:b1:alcohol-activation:sulfonate}
\emph{सल्फ़ोनेट एस्टर}\index{सल्फ़ोनेट एस्टर} $\mathrm{R{-}O{-}SO_2{-}R'}$
किसी ऐल्कोहॉल और सल्फ़ोनिल क्लोराइड $\mathrm{R'SO_2Cl}$ से किसी
क्षारक (पिरिडीन) की उपस्थिति में बनता है। \emph{टोसिलेट}\index{टोसिलेट}, \ce{R-OTs},
4-मेथिलबेंज़ीनसल्फ़ोनिल क्लोराइड (टोसिल क्लोराइड) से आता है;
\emph{मेसिलेट}\index{मेसिलेट}, \ce{R-OMs}, मेथेनसल्फ़ोनिल क्लोराइड से।
सल्फ़ोनेट ऋणायन $\mathrm{R'SO_3^-}$, जो एक \omterm{def:b1:predominant-reaction:strong-acid}{प्रबल अम्ल} का संयुग्मी क्षारक है और
तीन ऑक्सीजनों पर अनुनाद द्वारा स्थायी होता है, एक उत्कृष्ट \omterm{def:b1:electronic-effects:nucleophile}{निर्गामी समूह} है।
\end{definition}

\begin{omfigure}
\setchemfig{atom sep=2em}
\schemestart
\chemfig{R-O-H}
\+
\chemfig{Cl-S(=[2]O)(=[6]O)-[:0]*6(-=-(-CH_3)=-=)}
\arrow{->[पिरिडीन]}
\chemfig{R-O-S(=[2]O)(=[6]O)-[:0]*6(-=-(-CH_3)=-=)}
\schemestop
\par\vspace{6pt}

{\small किसी ऐल्कोहॉल का टोसिलन। O--H आबंध O--S से प्रतिस्थापित होता है;
ऐल्कोहॉल के C--O आबंध को छुआ नहीं जाता। पिरिडीन बने HCl को ले लेता है।}
\end{omfigure}

\begin{proposition}[टोसिलन और प्रतिस्थापन के दौरान विन्यास]\label{prop:b1:alcohol-activation:tosylation-retention}
त्रिविमजनक कार्बन पर किसी ऐल्कोहॉल का टोसिलन उसका \omterm{def:b1:stereochemistry-in-depth:isomers}{विन्यास} बनाए रखता है
(धारण); \omterm{def:b1:alcohol-activation:sulfonate}{टोसिलेट} का बाद का SN2 प्रतिस्थापन उसे प्रतीपित कर देता है।
दोनों चरण मिलकर OH को एक \omterm{def:b1:electronic-effects:nucleophile}{नाभिकरागी} से प्रतीपन के साथ प्रतिस्थापित करते हैं।
\end{proposition}

\begin{proof}
टोसिलन में जो आबंध बदलते हैं वे O--H और S--Cl हैं: ऐल्कोहॉल का
ऑक्सीजन सल्फ़र पर आक्रमण करता है और उसका हाइड्रोजन क्षारक के पास जाता है।
त्रिविमजनक कार्बन से कोई आबंध नहीं टूटता, इसलिए उसके चारों ओर की व्यवस्था
वही रहती है। SN2 चरण में C--O आबंध टूटता है जबकि \omterm{def:b1:electronic-effects:nucleophile}{नाभिकरागी}
विपरीत ओर आबंध बनाता है (\cref{prop:b1:nucleophilic-substitution:sn2-inversion})।
\end{proof}

\begin{omfigure}
\begin{tikzpicture}[font=\small, >={Stealth}]
  \node (a) at (0,0) {\chemfig{C(-[:150]H_3C)(<[:210]H)(<:[:250]C_2H_5)-OH}};
  \node at (0,-1.4) {\cip{S}-ब्यूटेन-2-ऑल};
  \draw[->, thick] (1.3,0) -- (2.5,0) node[midway, above] {TsCl} node[midway, below, font=\scriptsize] {धारण};
  \node (b) at (3.9,0) {\chemfig{C(-[:150]H_3C)(<[:210]H)(<:[:250]C_2H_5)-OTs}};
  \node at (3.9,-1.4) {\cip{S} टोसिलेट};
  \draw[->, thick] (5.3,0) -- (6.6,0) node[midway, above] {\ce{CN-}} node[midway, below, font=\scriptsize] {SN2, प्रतीपन};
  \node (c) at (8.3,0) {\chemfig{[:0]NC-C(-[:30]CH_3)(<[:-30]H)(<:[:-70]C_2H_5)}};
  \node at (8.3,-1.4) {\cip{R} नाइट्राइल};
\end{tikzpicture}

{\small \cip{S}-ब्यूटेन-2-ऑल से 2-मेथिलब्यूटेननाइट्राइल तक: टोसिलन
\omterm{def:b1:stereochemistry-in-depth:isomers}{विन्यास} बनाए रखता है, सायनाइड के साथ SN2 चरण उसे प्रतीपित कर देता है। कुल मिलाकर:
OH, प्रतीपन के साथ, CN से प्रतिस्थापित।}
\end{omfigure}

यही लक्ष्य अम्ल में भी पाया जा सकता है: प्रोटॉनित OH समूह
जल के रूप में निकलता है, \ce{R-OH2+ -> R+ + H2O} (SN1, E1), या आक्रमण के अंतर्गत (SN2)। पर अम्ल
\omterm{def:b1:electronic-effects:nucleophile}{नाभिकरागी} को उन्हीं तक सीमित कर देता है जो उसमें टिक सकें --- हैलाइड आयन, जल,
ऐल्कोहॉल --- और \omterm{def:b1:electronic-effects:intermediates}{कार्बधनायनों} तथा उनकी पार्श्व अभिक्रियाओं का द्वार खोल देता है।

\section{हैलोऐल्केनों में रूपांतरण}

\begin{proposition}[ऐल्कोहॉल और हाइड्रोजन हैलाइड]\label{prop:b1:alcohol-activation:hx-by-class}
तृतीयक ऐल्कोहॉल कमरे के ताप पर सांद्र हाइड्रोक्लोरिक
अम्ल के साथ हैलोऐल्केन देते हैं (SN1); प्राथमिक ऐल्कोहॉलों को हाइड्रोब्रोमिक या
हाइड्रोआयोडिक अम्ल और गरम करना चाहिए (प्रोटॉनित ऐल्कोहॉल पर SN2); द्वितीयक
ऐल्कोहॉल मध्यवर्ती दरों पर अभिक्रिया करते हैं, प्रायः दोनों क्रियाविधियों से।
\end{proposition}

\begin{proof}
प्रोटॉनन OH को \ce{-OH2+} में बदल देता है। तृतीयक ऐल्कोहॉल के लिए जल
आसानी से निकलता है, जिससे एक स्थायी \omterm{def:b1:electronic-effects:intermediates}{कार्बधनायन} बनता है जिसे हैलाइड ग्रहण कर लेता है (SN1): साधारण
\omterm{def:b1:electronic-effects:nucleophile}{नाभिकरागी} \ce{Cl-} भी पर्याप्त है। प्राथमिक \omterm{def:b1:electronic-effects:intermediates}{कार्बधनायन} बहुत
अस्थायी है; हैलाइड को जल को पीछे से धकेलकर बाहर निकालना होता है (SN2), जिसके लिए
अच्छा \omterm{def:b1:electronic-effects:nucleophile}{नाभिकरागी} (\ce{Br-}, \ce{I-}, प्रोटिक माध्यम में \ce{Cl-} से बेहतर,
\cref{prop:b1:electronic-effects:nucleophilicity-basicity}) और ऊष्मा चाहिए।
\end{proof}

\begin{example}[एक वर्गीकरण परीक्षण]\label{ex:b1:alcohol-activation:lucas}
सांद्र हाइड्रोक्लोरिक अम्ल और ज़िंक क्लोराइड (एक \omterm{def:b1:lewis-resonance-vsepr:lewis-acid}{लुइस
अम्ल} जो OH को निकलने में सहायता करता है) का मिश्रण तृतीयक ऐल्कोहॉल को तुरंत धुँधला कर देता है ---
अविलेय क्लोराइड अलग हो जाता है --- द्वितीयक ऐल्कोहॉल को कुछ मिनटों में,
और प्राथमिक ऐल्कोहॉल को कमरे के ताप पर शायद ही: \omterm{def:b1:electronic-effects:intermediates}{कार्बधनायन}
स्थायित्व का क्रम एक परखनली में दृश्य हो उठता है।
\end{example}

\section{अम्ल में निर्जलीकरण और ईथर का बनना}

\begin{definition}[ऐल्कोहॉल का निर्जलीकरण]\label{def:b1:alcohol-activation:dehydration}
किसी ऐल्कोहॉल का \emph{निर्जलीकरण}\index{निर्जलीकरण} जल का
विलोपन है, जो किसी \omterm{def:b1:predominant-reaction:strong-acid}{प्रबल अम्ल} (सल्फ़्यूरिक या फ़ॉस्फ़ोरिक
अम्ल) और ऊष्मा से उत्प्रेरित होता है और एक ऐल्कीन देता है: $\mathrm{R_2CH{-}CR'_2OH} \to \mathrm{R_2C{=}CR'_2} + \ce{H2O}$।
\end{definition}

\begin{omfigure}
\setchemfig{atom sep=2.1em}
\schemestart
\chemfig{C(-[2]CH_3)(-[:210]H_3C)(-[:-30]CH_2CH_3)-OH}
\arrow{->[\ce{H+}]}
\chemfig{C(-[2]CH_3)(-[:210]H_3C)(-[:-30]CH_2CH_3)-OH_2^{+}}
\arrow{->[\ce{-H2O}]}
\chemfig{C^{+}(-[2]CH_3)(-[:210]H_3C)(-[:-30]CH_2CH_3)}
\schemestop

\medskip
\schemestart
\chemfig{C^{+}(-[2]CH_3)(-[:210]H_3C)(-[:-30]CH_2CH_3)}
\arrow{->[\ce{-H+}][(मुख्य)]}
\chemfig{C(-[2]CH_3)(-[:210]H_3C)=[:-30]CHCH_3}
\schemestop

\medskip
{\small 2-मेथिलब्यूटेन-2-ऑल का अम्ल-उत्प्रेरित \omterm{def:b1:alcohol-activation:dehydration}{निर्जलीकरण} (E1): प्रोटॉनन,
तृतीयक \omterm{def:b1:electronic-effects:intermediates}{कार्बधनायन} तक जल की हानि, एक प्रोटॉन की हानि।
\ce{CH2} से प्रोटॉन हटाने पर त्रिप्रतिस्थापित 2-मेथिलब्यूट-2-ईन मिलती है
(ज़ैत्सेव, मुख्य); किसी मेथिल समूह से हटाने पर 2-मेथिलब्यूट-1-ईन (गौण)।}
\end{omfigure}

\begin{proposition}[ईथर या ऐल्कीन]\label{prop:b1:alcohol-activation:ether-vs-alkene}
सल्फ़्यूरिक अम्ल के साथ गरम किया गया प्राथमिक ऐल्कोहॉल मध्यम ताप पर मुख्यतः एक सममित
ईथर देता है (एक अणु SN2 द्वारा दूसरे का प्रतिस्थापन करता है)
और उच्चतर ताप पर मुख्यतः एक ऐल्कीन; द्वितीयक और तृतीयक
ऐल्कोहॉल ऐल्कीन देते हैं।
\end{proposition}

\begin{proof}
प्राथमिक ऐल्कोहॉल के लिए, प्रोटॉनित ऐल्कोहॉल पर या तो कार्बन पर
ऐल्कोहॉल के दूसरे अणु द्वारा आक्रमण हो सकता है (SN2, जो प्रोटॉन की हानि के बाद \ce{R-O-R}
देता है) या किसी $\beta$ हाइड्रोजन पर (E2)। विलोपन जितने अणु खपाता है उससे अधिक
बनाता है और ताप बढ़ने के साथ इसे वरीयता मिलती है
(\cref{prop:b1:elimination:sn-vs-e}); ताप कम होने पर प्रतिस्थापन
प्रभावी होता है। द्वितीयक और तृतीयक ऐल्कोहॉल आयनित होकर \omterm{def:b1:electronic-effects:intermediates}{कार्बधनायन} बनाते हैं, जो
किसी दुर्बल \omterm{def:b1:electronic-effects:nucleophile}{नाभिकरागी} द्वारा ग्रहण किए जाने की तुलना में प्रोटॉन अधिक आसानी से खो देते हैं, और
गरम करने पर यह और अधिक होता है; बना हुआ जल भी वाष्पशील ऐल्कीन के
आसवन से हट जाता है।
\end{proof}

\section{अभ्यास}

\begin{exercise}[$\star$]\label{exo:b1:alcohol-activation:1}
सोडियम के साथ, सोडियम हाइड्राइड के साथ, और सोडियम हाइड्रॉक्साइड के साथ एथेनॉल की
अभिक्रियाएँ लिखिए; अंतिम का स्थिरांक परिकलित कीजिए (एथेनॉल का $\mathrm{p}K_a$
15.9)।
\end{exercise}

\begin{exercise}[$\star$]\label{exo:b1:alcohol-activation:2}
इनके उत्पाद बताइए: 1-ब्रोमोप्रोपेन के साथ सोडियम मेथॉक्साइड; आयोडोमेथेन के साथ सोडियम
एथॉक्साइड; ब्रोमोएथेन के साथ सोडियम फ़ीनॉक्साइड।
\end{exercise}

\begin{exercise}[$\star$]\label{exo:b1:alcohol-activation:3}
प्रोपेन-1-ऑल का टोसिलन और सोडियम आयोडाइड के साथ \omterm{def:b1:alcohol-activation:sulfonate}{टोसिलेट} की अभिक्रिया
लिखिए। सोडियम आयोडाइड के साथ ऐल्कोहॉल की प्रत्यक्ष अभिक्रिया की तुलना में
इसका क्या लाभ है?
\end{exercise}

\begin{exercise}[$\star$]\label{exo:b1:alcohol-activation:4}
ब्यूटेन-2-ऑल और 2-मेथिलपेंटेन-2-ऑल के अम्लीय \omterm{def:b1:alcohol-activation:dehydration}{निर्जलीकरण} से बनने वाली ऐल्कीन बताइए,
और उनमें मुख्य कौन-सी है।
\end{exercise}

\begin{exercise}[$\star\star$]\label{exo:b1:alcohol-activation:5}
1-एथॉक्सी-2-मेथिलप्रोपेन का एक विलियमसन संश्लेषण प्रस्तावित कीजिए और
भागीदारों के अपने चयन की व्याख्या कीजिए। दूसरा चयन ऐल्कीन क्यों देता?
\end{exercise}

\begin{exercise}[$\star\star$]\label{exo:b1:alcohol-activation:6}
\cip{R}-ऑक्टेन-2-ऑल को उसके \omterm{def:b1:alcohol-activation:sulfonate}{टोसिलेट} में बदला जाता है, फिर
एथेनॉल में सोडियम एथॉक्साइड से उपचारित किया जाता है। ईथर का \omterm{def:b1:stereochemistry-in-depth:isomers}{विन्यास} बताइए। उसी
ऐल्कोहॉल को सल्फ़्यूरिक अम्ल के साथ गरम करने पर क्या मिलता?
\end{exercise}

\begin{exercise}[$\star\star$]\label{exo:b1:alcohol-activation:7}
सांद्र हाइड्रोक्लोरिक अम्ल के साथ 2-मेथिलप्रोपेन-2-ऑल की अभिक्रिया की
क्रियाविधि लिखिए और समझाइए कि कमरे के ताप पर अभिक्रिया तीव्र क्यों
है।
\end{exercise}

\begin{exercise}[$\star\star$]\label{exo:b1:alcohol-activation:8}
सल्फ़्यूरिक अम्ल में एथेनॉल से एथॉक्सीएथेन बनने की क्रियाविधि, और
प्रतिस्पर्धी विलोपन, लिखिए।
\end{exercise}

\begin{exercise}[$\star\star$]\label{exo:b1:alcohol-activation:9}
इस अध्याय के वर्गीकरण परीक्षण में ब्यूटेन-1-ऑल, ब्यूटेन-2-ऑल और
2-मेथिलप्रोपेन-2-ऑल में से कौन-सा ऐल्कोहॉल सबसे पहले धुँधला होता है? उत्पाद
लिखिए।
\end{exercise}

\begin{exercise}[$\star\star\star$]\label{exo:b1:alcohol-activation:10}
3,3-डाइमेथिलब्यूटेन-2-ऑल, अम्ल में गरम करने पर, मुख्यतः
2,3-डाइमेथिलब्यूट-2-ईन देता है, जिसका कार्बन कंकाल
ऐल्कोहॉल के कंकाल से भिन्न है। \omterm{def:b1:electronic-effects:intermediates}{कार्बधनायन} में एक मेथिल समूह के 1,2-विस्थापन वाली
एक क्रियाविधि प्रस्तावित कीजिए, और विस्थापन का औचित्य दीजिए।
\end{exercise}

\begin{exercise}[$\star\star\star$]\label{exo:b1:alcohol-activation:11}
\cip{R}-ब्यूटेन-2-ऑल से \cip{S}-2-मेथॉक्सीब्यूटेन का एक संश्लेषण अभिकल्पित कीजिए,
और एक अन्य \cip{S}-ब्यूटेन-2-ऑल से, दोनों में \omterm{def:b1:stereochemistry-in-depth:isomers}{विन्यास} को
नियंत्रण में रखते हुए।
\end{exercise}

\begin{exercise}[$\star\star\star$]\label{exo:b1:alcohol-activation:12}
दिखाइए कि साम्य \ce{2CH3CH2OH <=> (CH3CH2)2O + H2O} को
जल को लगातार हटाकर ईथर की ओर धकेला जा सकता है, \omterm{def:b1:extent-q-and-k:quotient}{अभिक्रिया
भागफल} (\cref{ch:b1:extent-q-and-k}) का उपयोग करते हुए।
\end{exercise}

\section{समस्या: एक ईथर तक दो मार्ग}

\begin{problem}[{सप्ताहांत समस्या --- मेथिल
टर्ट-ब्यूटिल ईथर तक दो विलियमसन मार्ग, प्रकाशतः सक्रिय ऐल्कोहॉल से
त्रिविम-रासायनिक रूप से नियंत्रित एक ईथर, एक तृतीयक ऐल्कोहॉल का निर्जलीकरण, और
प्राप्त ईथर का द्रव्यमान}]
\label{pb:b1:alcohol-activation:1}
2-मेथॉक्सी-2-मेथिलप्रोपेन (मेथिल टर्ट-ब्यूटिल ईथर, MTBE) ईंधन योज्य के रूप में
बड़े पैमाने पर बनाया जाता था। मोलर द्रव्यमान (\unit{g/mol}): \ce{C4H10O}
74.0, \ce{C5H12O} 88.0, \ce{CH3I} 141.9।

\medskip
\noindent\textbf{भाग I --- विलियमसन संश्लेषण द्वारा MTBE।}
\begin{enumerate}
  \item भागीदारों (\omterm{def:b1:alcohol-activation:alkoxide}{ऐल्कॉक्साइड} और हैलाइड) के वे दो जोड़े लिखिए जो
        MTBE दे सकते हैं।
  \item कौन-सा जोड़ा विफल होता है? उसके बजाय होने वाली अभिक्रिया लिखिए।
  \item सफल अभिक्रिया और उसकी क्रियाविधि लिखिए।
  \item 2-मेथिलप्रोपेन-2-ऑल से सोडियम टर्ट-ब्यूटॉक्साइड कैसे बनाया जाता है?
        सोडियम हाइड्रॉक्साइड से क्यों नहीं?
  \item अभिक्रिया शुष्क ऐल्कोहॉल या किसी \omterm{def:b1:intermolecular-forces-solvents:solvent-classes}{अप्रोटिक विलायक} में क्यों की जाती है?
  \item क्या उत्पाद \omterm{def:b1:stereochemistry-in-depth:stereocentre}{किरैल} है?
\end{enumerate}

\medskip
\noindent\textbf{भाग II --- एक प्रकाशतः सक्रिय ईथर।}
\begin{enumerate}[resume]
  \item \cip{R}-ब्यूटेन-2-ऑल को उसके \omterm{def:b1:alcohol-activation:sulfonate}{टोसिलेट} में बदला जाता है। अभिक्रिया
        लिखिए और \omterm{def:b1:alcohol-activation:sulfonate}{टोसिलेट} का \omterm{def:b1:stereochemistry-in-depth:isomers}{विन्यास} बताइए।
  \item \omterm{def:b1:alcohol-activation:sulfonate}{टोसिलेट} सोडियम मेथॉक्साइड से अभिक्रिया करता है। उत्पाद और
        उसका \omterm{def:b1:stereochemistry-in-depth:isomers}{विन्यास} लिखिए।
  \item उसी मार्ग से ऐल्कोहॉल का कौन-सा \omterm{def:b1:stereochemistry-in-depth:isomers}{विन्यास} \cip{R}-2-मेथॉक्सीब्यूटेन
        देता है?
  \item एक अन्य मार्ग: \cip{R}-ब्यूटेन-2-ऑल का \omterm{def:b1:alcohol-activation:alkoxide}{ऐल्कॉक्साइड}
        आयोडोमेथेन के साथ। उत्पाद का \omterm{def:b1:stereochemistry-in-depth:isomers}{विन्यास}? क्यों?
  \item प्रश्न 8 में कौन-सी पार्श्व अभिक्रिया प्रतिस्पर्धा करती है, और यहाँ वह
        सीमित क्यों है?
  \item \cip{R}-ब्यूटेन-2-ऑल को मेथेनॉल के साथ सल्फ़्यूरिक
        अम्ल में गरम करने से, यदि कोई ईथर बने भी, तो लगभग रेसेमिक क्यों होगा?
\end{enumerate}

\medskip
\noindent\textbf{भाग III --- एक तृतीयक ऐल्कोहॉल का \omterm{def:b1:alcohol-activation:dehydration}{निर्जलीकरण}।}
\begin{enumerate}[resume]
  \item 2-मेथिलब्यूटेन-2-ऑल के अम्लीय \omterm{def:b1:alcohol-activation:dehydration}{निर्जलीकरण} की
        क्रियाविधि लिखिए।
  \item दोनों ऐल्कीन बताइए और मुख्य का पूर्वानुमान कीजिए।
  \item इस ऐल्कोहॉल से कोई ईथर क्यों नहीं बनता?
  \item 2-मेथिलब्यूट-2-ईन को बनते ही आसवित करके क्यों निकाल लिया जाता है?
  \item E1 चरण का \omterm{def:b1:elementary-steps:energy-profile}{ऊर्जा प्रोफ़ाइल} खींचिए और
        \omterm{def:b1:elementary-steps:rds}{दर-निर्धारक चरण} चिह्नित कीजिए।
  \item क्या किसी प्राथमिक ऐल्कोहॉल का \omterm{def:b1:alcohol-activation:dehydration}{निर्जलीकरण} इसी क्रियाविधि से हो सकता है?
\end{enumerate}

\medskip
\noindent\textbf{भाग IV --- द्रव्यमान।}
\begin{enumerate}[resume]
  \item \qty{10.0}{g} 2-मेथिलप्रोपेन-2-ऑल को
        \omterm{def:b1:alcohol-activation:alkoxide}{ऐल्कॉक्साइड} में बदला जाता है। मात्रा परिकलित कीजिए।
  \item \qty{10}{\percent} आधिक्य के लिए आयोडोमेथेन का कितना द्रव्यमान चाहिए?
  \item MTBE का सैद्धांतिक द्रव्यमान परिकलित कीजिए।
  \item \qty{75}{\percent} \omterm{def:b1:extent-q-and-k:fractional-extent}{लब्धि} के लिए प्राप्त द्रव्यमान परिकलित कीजिए।
\end{enumerate}
\end{problem}
